\subsection{Normal algebraic extensions and the remaining subgroup boundary}

\begin{claim}{B.META.17}{Proved}{An abelian normal quotient supplies the arithmetic inputs for reconstruction}
Keep $D$, $M_0$, $\Omega$ and
$I=I_{M_0}\subset G_D$ as in Claim~B.META.15.
Let $H$ be a nontrivial closed normal subgroup of
$G_{M_0}$, and let $B$ be a closed normal subgroup
of $I$ such that
\[
 H\subsetneq B\subset I\cap\ker\chi_p,
                         \qquad [B,B]\subset H.
                                                        \tag{META.NG1}
\]
Put $K=\Omega^H$ and $A=\Omega^B$.
Then every finite extension of $K$ satisfies full
closed-point-preserving relative etale reconstruction
for finite-type schemes localized at universal
homeomorphisms, with the structural isomorphism
to the base retained.

The arithmetic assertion needed here is the following
one-sided description of the stable rational Kummer
kernel: for every finite $E/K$ and every $a\in E^\times$,
\[
 \kappa_E(a)=0\quad\Longrightarrow\quad
  g(a)/a\in\mu_\infty\ (g\in B),
 \quad a^m\in A\text{ for some }m\ge1.
                                                        \tag{META.NG2}
\]
No equality between this kernel and all radicals
of $A$ is asserted.
\end{claim}
\begin{proof}
The inertia group fixes all prime-to-$p$ roots
of unity, and $B$ also fixes all $p$-power roots
by its cyclotomic condition.  Hence
$\mu_\infty\subset A\subset K$.
All these fields are algebraic extensions of
the henselian field $D^h$, with the fixed
mixed-characteristic valuation.

Suppose that $a\in E^\times$ has zero rational
Kummer class.  Choose a finite Galois extension
$D_1/D$ inside $\Omega$ containing $a$ and
generators of $E/K$.  Let $M_1$ be its finite
local closure, and replace $E$ by $KD_1$.
Set
\[
 J=I\cap G_{D_1}=I_{M_1},\quad H_1=H\cap J,
 \quad\Gamma=J/H_1,\quad C=(B\cap J)/H_1.
\]
The local class of $a$ in $H^1(J,V)$, where
$V=\mathbb Q_p(1)$, vanishes on $H_1$ and
inflates from a unique $\theta\in H^1(\Gamma,V)$.
The group $C$ is abelian and acts trivially on $V$.
For any $g\in B$, Galoisness of $D_1/D$ makes
$g$ normalize $J$; normality of $H$ and $B$
makes the resulting action preserve $\Gamma$
and $C$.
It acts trivially on $C$, because $[B,B]\subset H$,
and on $\Gamma/C$, because $B$ is normal in $I$.
It acts trivially on $V$ by $\chi_p(g)=1$.

Inflation--restriction gives
\[
 0\longrightarrow H^1(\Gamma/C,V)
 \longrightarrow H^1(\Gamma,V)
 \longrightarrow H^1(C,V)^{\Gamma/C}.
                                                        \tag{META.NG3}
\]
The action of $g$ is the identity on both outer
terms.  Thus $(g-1)^2=0$ on the middle term.
Since $a\in D_1$ has a finite $g$-orbit, there
is $s\ge1$ with $g^s(a)=a$.
Naturality of Kummer theory and injectivity of
inflation give $g^s\theta=\theta$.
The unipotence identity now yields
\[
               0=(g^s-1)\theta=s(g-1)\theta,
\]
so $g\theta=\theta$ in characteristic zero.
This uses a finite orbit only for the specific
Kummer vector, not for all of $H^1(\Gamma,V)$.
We obtain $\kappa_J(g(a)/a)=0$.

The elements $a$ and $g(a)$ belong to the one
finite local field $M_1$, since $g$ preserves
$D_1$ and its valuation closure.
Restriction from $H^1(G_{M_1},V)$ to
$H^1(I_{M_1},V)$ is injective: the coefficient
module has no inertia invariants.
The rational Kummer kernel of $M_1$ is its
finite group of roots of unity.
Indeed valuation detects the uniformizer part,
and the principal-unit logarithm detects an
open subgroup of units modulo its finite
torsion; the residue unit group is finite.
Hence $g(a)/a\in\mu(M_1)$ for every $g\in B$.
One integer annihilating $\mu(M_1)$ makes
$a^m$ fixed by all of $B$, proving (META.NG2).
The argument included every finite $E/K$.

Choose $b\in K\setminus A$, which exists by
$H\subsetneq B$.
At least one of $b$ and $1-b$ has nonzero
Kummer class over $K$.
Otherwise (META.NG2) would give, for every $g\in B$,
\[
        g(b)=\zeta_g b,\qquad
        g(1-b)=\eta_g(1-b),
                 \quad\zeta_g,\eta_g\in\mu_\infty\subset A.
\]
If some pair is unequal, its sum expresses
$b=(1-\eta_g)/(\zeta_g-\eta_g)\in A$.
If every pair is equal, the sum makes each
pair one and all $g$ fix $b$, again a contradiction.
This nonzero class remains nonzero over every
finite extension by rational restriction and
corestriction.  In particular all the fields
used below are non-separably-closed.

For a degree $d\ge1$ use
$\mathcal G_0=\{u+iv:0\le u,v\le4d\}$ and
$\mathcal G_1=b+\mathcal G_0$.
Suppose a quotient $R=P/Q$ of monic degree-$d$
polynomials over $\Omega$ has nonzero grid
values and all those values are zero Kummer
classes after finite extensions of $K$.
For every $g\in B$, (META.NG2) gives
$g(R(a))/R(a)\in\mu_\infty$.
Since $g$ fixes the first grid, write
$gR/R=U/V$ with monic polynomials of equal
degree at most $2d$ before cancellation.
Under an embedding $\Omega\hookrightarrow\mathbb C$,
the polynomial
\[
 U(u+iv)\overline U(u-iv)-V(u+iv)\overline V(u-iv)
\]
vanishes on the whole grid and has degree at
most $4d$ separately in $u,v$.
Successive interpolation makes it zero.
The invertible change $z=u+iv,w=u-iv$ gives
$U(z)\overline U(w)=V(z)\overline V(w)$.
Fixing an allowed $w$ shows proportionality,
and monicity gives $gR=R$.
Repeating this on $b+\mathcal G_0$ gives
$gR(gb+Z)=R(b+Z)$.
Choose $g$ moving $b$.
Then $R$ has the nonzero additive period $gb-b$.
A finite pole would have infinitely many
translates; a nonconstant polynomial has a
nonzero leading finite difference.
Thus $R$ is constant, and its value at infinity
is one.  We have proved the uniform two-grid test.

We now verify the remaining inputs of the
norm proof, rather than treating (META.NG2)
alone as reconstruction.
For a smooth geometrically integral affine
source and a pinned map to $\mathbb A^1$, fix
any finite Noether projection and the two
grids just constructed.
As in (META.RL4), choose a principal open
avoiding the branch locus and the proper
closed images of the grid fibres.
The geometric unit calculation, whose Picard
group is zero, gives norm classes
\[
                 \xi_a=\beta_a+
                       \sum_j\lambda_{aj}\kappa(h_j),
                  \qquad\lambda_{aj}\in\mathbb Q_p.
\]
Include all scheme and boundary coefficients,
grid centres, one algebraic normalizing parameter
and its actual norm values in a finite Galois
$D_1/D$.  Evaluation at that parameter identifies
each $\beta_a$ with an explicit finite linear
combination of Kummer classes from $D_1$.
With $E_1=KD_1$ and $M_1=\overline{D_1}$
in the local topology, the group
$N_1=G_{E_1}=H\cap G_{D_1}$ is normal in
$G_{M_1}$ and fixes $\mathbb Q_p(1)$.
All normalized classes are $G_{M_1}$-invariant.
Comparison of their evaluations at $t$ and
$\sigma t$ puts every grid value of
$P_{\sigma t}/\sigma P_t$ in the stable kernel.
The two-grid test gives
$P_{\sigma t}=\sigma P_t$ uniformly in
$\sigma\in G_{M_1}$ and the parameter.

The nonzero Kummer class constructed above
restricts nontrivially to $N_1$.
Dividing its integral image $p^r\mathbb Z_p$
by taking an actual $p^r$-th root in $E_1$
gives a Kummer homomorphism onto $\mathbb Z_p$.
After adjoining its algebraic argument and
$\mu_{p^s}$ to the finite local field, it
produces a cyclic $L/M_1$ of any prescribed
degree $q=p^s$, with surjective image of $N_1$.
These enlargements lie in the local fixed
field of $N_1$, so normality persists.
Claim~B.META.15 puts $M_1\cap\Omega$ in
the henselian $E_1$, identifies
$L\cap\Omega$ as a finite extension of it,
and supplies an $M_1$-basis of $L$ in that
algebraic part.

If a class over $L$ vanishes on
$N_1\cap G_L$, it inflates from the quotient.
The group $N_1$ acts trivially on that quotient
and on the coefficient module.
Its local logarithm is therefore in $M_1$,
because $N_1$ surjects onto $\operatorname{Gal}(L/M_1)$.
Actual point maps are valuation-continuous
over all finite extensions: the qcs etale-site
and relative-point comparison of CHW gives
etale-open continuity, and JTWY Corollary~6.13
applies to these non-separably-closed henselian
fields \citep{AN_CHW2024,AN_JTWY2024}.

All inputs of the norm proof in
Claim~B.META.14 have now been verified over
$\Omega$, with the local comparison of
Claim~B.META.15.
To specify the use, take $q>d+1$.
Continuity and minimum conjugate gcd degree
give a constant-degree neighbourhood on the
dense algebraic parameter set.
Exact local logarithms and $2e$-centre
interpolation recover the reduced ratios.
Their extended coordinate derivations and
cocycle give the root divisor
$\tau^jB_0-B_0+C_j$.
Constant roots are algebraic over $D$ by
one actual specialization and are excluded
by proper closed actual fibres.
A negative coefficient of $B_0$ would give
more than $d$ distinct denominator roots;
specialization off the nonzero analytic
subresultants contradicts the actual degree.
The resulting effective divisor gives the
analytic monic denominator $Q$ in one block.
The exact logarithmic equations satisfy
B.135, so $Q\in L(t)[Z]$; interpolation of
its values in $\Omega$ gives $Q\in\Omega(t)[Z]$.
Residues are integral, and one actual
normalization gives an arithmetic constant
factor $C$ with $P_t=Q_tC$ by the proved
two-grid test.
Proper closed fibres force $C=1$.
This recovers all finite-projection norms.
The fully proved graph-degree argument in
Claim~B.META.14 makes their graph a single
rational function.  Valuation continuity
on algebraic transverse curves removes
poles and exceptional point values.
It is a regular function over the original
base by Galois descent.
Geometric components are handled by one
finite extension and the same uniqueness.
Finally CHW Theorem~4.18 and Corollary~2.32
give full reconstruction in the stated
relative category \citep{AN_CHW2024}.
Every step starts with any finite $E/K$,
as required.
\end{proof}

\begin{claim}{B.META.18}{Proved}{Every proper Galois local algebraic extension reconstructs}
Keep the data of Claim~B.META.15.
Let $1\ne N\triangleleft G_{M_0}$ be any
closed normal subgroup.
Then $F=\Omega^N$, every finite extension of
$F$ in $\Omega$, and every algebraic subfield
of any of those fields satisfy the reconstruction
property of Claim~B.META.17.
Equivalently, every proper Galois algebraic
extension of $D^h=M_0\cap\Omega$ reconstructs,
with the same finite-extension and algebraic
descent assertions.

More generally let $D^h\subset F\subset\Omega$.
If for some finite local extension $M_1/M_0$
the closed subgroup
\[
 \operatorname{Core}_{G_{M_1}}(G_F\cap G_{M_1})
   =\bigcap_{g\in G_{M_1}}g(G_F\cap G_{M_1})g^{-1}
                                                        \tag{META.NG4}
\]
is nontrivial, then $F$ and every finite
extension of $F$ reconstruct.
Thus within these algebraic local subfields
the remaining cases have trivial normal
core in every open local absolute Galois
group.
\end{claim}
\begin{proof}
We first show that every nontrivial closed
normal subgroup of a finite-local-field
absolute Galois group is nonabelian.
Let $P=P_{M_0}$ be its wild inertia and
suppose $N$ is abelian.
Then $A=N\cap P$ is abelian and normal in $P$.
If $a\in A$ is nonidentity, choose two
noncommuting elements $s,t\in P$ and let
$P_0=\overline{\langle a,s,t\rangle}$.
Wild inertia is free pro-$p$ of infinite
rank, and its finitely generated closed
subgroup $P_0$ is free pro-$p$.
The nontrivial abelian normal subgroup
$A\cap P_0$ is procyclic, because it lies
in the procyclic centralizer of $a$ in
$P_0$.
The normalizer of a nonidentity procyclic
subgroup in a finitely generated free
pro-$p$ group is procyclic as well.
Thus normality would make $P_0$ procyclic,
contrary to the choice of $s,t$.
These exact free pro-$p$ statements are
Iwasawa Theorem~2(ii) and Kochloukova--Zalesskii
Theorem~5.1, Case~I, opened and used with
the same finite-generator reduction in
Claim~B.META.1
\citep{AN_Iwasawa1955,AN_KZ2011}.
We obtain $A=1$.
Normality of both $N$ and $P$ gives
$[N,P]\subset N\cap P=1$.
The local centralizer theorem
$C_{G_{M_0}}(P)=1$, proved at every prime
in Claim~B.META.1, now gives $N=1$.
This contradiction proves the assertion.
The argument applies to every finite
extension of $M_0$ as well.

Return to a nontrivial normal $N$ and put
$B=N'=\overline{[N,N]}$.
It is nontrivial by what we just proved,
and it is normal in $G_{M_0}$.
The quotient $G_{M_0}/I$ and the cyclotomic
image are abelian, so
\[
                         B\subset I\cap\ker\chi_p.
\]
We do not assert that $B$ is pro-$p$.
Every finite quotient of $I$ is an extension
of a finite $p$-group by a cyclic prime-to-$p$
group, hence solvable.
It follows that $\bigcap_r I^{(r)}=1$:
a nonidentity element has a nonidentity
image in some finite solvable quotient,
whose derived series terminates.
There is consequently a least integer
$r\ge1$ with $B\not\subset I^{(r)}$.
Minimality gives $B\subset I^{(r-1)}$.
Set
\[
                              H=B\cap I^{(r)}.
\]
This is a closed normal subgroup of $G_{M_0}$,
proper in $B$.
The commutator inclusion
$[B,B]\subset [I^{(r-1)},I^{(r-1)}]\subset I^{(r)}$
makes $B/H$ abelian and nontrivial.
Moreover $H\ne1$:
the nontrivial normal subgroup $B$ is
nonabelian by the first paragraph, so
$1\ne B'\subset H$.
Every hypothesis of Claim~B.META.17 has
therefore been proved.
It reconstructs $K=\Omega^H$ and all its
finite extensions.
Since $H\subset N$, the field $K$ contains
$F=\Omega^N$, and $K/F$ is algebraic.
The explicit algebraic-descent argument
in Claim~B.META.16 now proves reconstruction
over $F$.
For any finite $F'/F$, the compositum
$KF'/K$ is finite; reconstruction there
descends to $F'$.
The same argument gives the stated
algebraic subfield assertion.

For (META.NG4), write its nontrivial core
as $C$.  It is a closed normal subgroup
of $G_{M_1}$.
By Claim~B.META.15 one may choose a finite
extension $D_1/D$ inside $\Omega$ with
closure $M_1$; its henselization is
$M_1\cap\Omega$.
Apply the first assertion with this base
and $C$.
The reconstructed proper Galois field
$\Omega^C$ contains $F$, since $C\subset G_F$.
Its extension of $F$ is algebraic.
Descent, and descent after every finite
compositum, prove the assertion.

Finally every henselian valued subfield
$F\subset\overline{\mathbb Q}_p$ of
finite transcendence degree contains the
henselization of a finitely generated
$D\subset F$ containing a transcendence
basis.  Then $F/D^h$ is algebraic and
the criterion just proved applies whenever
its indicated core is nontrivial.
No assertion for core-free subgroups is
hidden in this last reduction.
\end{proof}

\paragraph{Remaining scope.}
The normality argument is a new sufficient
criterion within the algebraic closure of
a finitely generated subfield of
$\overline{\mathbb Q}_p$.
It does not replace a closed subgroup by
its normal closure or its normal core
without checking that the core is nontrivial.
Nor does it embed arbitrary henselian
mixed-characteristic fields into
$\overline{\mathbb Q}_p$.
In particular fields with more general
residue fields or value groups, and the
remaining center-free core-free subgroup
cases, require separate arguments.
No priority or external human verification
is claimed.
